We contribute a reusable message-level functional analysis protocol for communication in two-agent LLM systems, evaluated across two primary backbone families, three reasoning domains, and 15 configurations, with 10-seed replication for the flagship settings. In the tested settings, function type is associated with pruning safety more reliably than surface content alone. Perturbation-pruning dissociation supports this finding: messages most disruptive to remove individually are safely prunable in aggregate, and acknowledgment-targeted pruning eliminates 36--56\% of inference calls at non-inferior accuracy ($\delta{=}{-}3$\,pp). The protocol also reveals backbone-dependent mechanisms not visible to content-level analysis: structural redundancy for Qwen and error-channel elimination for Llama. These results hold for the tested flagship configurations; boundary conditions (MATH, where acknowledgments carry substantive content; Llama under stochastic decoding with heuristic detection) and topology variants beyond two-agent teams remain underexplored.

We conjecture that aggregate removal dissolves self-reinforcing acknowledgment chains (Appendix~\ref{app:cascade}); this hypothesis requires prospective validation. All experiments use two-agent topologies under greedy decoding, and heuristic detectors require per-backbone calibration. A learned DistilBERT classifier bridges the oracle-to-deployment gap on the Qwen backbone; Llama deployment testing remains future work.

Because the protocol requires only a taxonomy-aware LLM and constrained decoding, practitioners can apply it to new backbones, domains, or multi-agent topologies without task-specific training data. Future work should extend to larger agent teams and verify this conjecture with prospective designs.
